\documentclass[10pt,a4paper]{amsart}
\usepackage[margin=19mm]{geometry}
\usepackage{amsmath,amssymb,mathtools}
\usepackage[scr=rsfs,cal=euler]{mathalfa}
\usepackage{xfrac}
\usepackage[unicode,hidelinks,bookmarks=false]{hyperref}
\newcommand{\ie}{i.e.\ }
\newcommand{\cf}{cf.\ }
\newcommand{\resp}{respectively\ }
\newcommand{\Subsection}{\subsection}
\providecommand{\nobreakdash}{\nobreak\hbox{-}}
\setlength{\emergencystretch}{2em}
\multlinegap0pt
\usepackage{bm}
\usepackage{bbm}
\usepackage{dsfont}
\usepackage{xspace}
\usepackage{cancel}
\usepackage{mathtools}
\usepackage{amsthm}
\makeatletter
\@ifundefined{theorem}{\newtheorem{theorem}{Theorem}}{}
\@ifundefined{claim}{\newtheorem{claim}[theorem]{Claim}}{}
\makeatother
\newcommand{\Ff}{\mathcal{G}}
\newcommand{\R}{{\mathbb R}} \newcommand{\Q}{{\mathbb Q}} \newcommand{\Z}{{\mathbb Z}} \newcommand{\N}{{\mathbb N}}
\newcommand{\Zer}{{\rm Zer}} 
\title[Stability of shifts, interpolation, and crystalline measures]{Stability of shifts, interpolation, and crystalline measures}
\author{Alexander Ulanovskii, Ilya Zlotnikov}
\date{Source: arXiv:2410.08682v2 (CC BY 4.0)}
\begin{document}
\maketitle

\noindent\emph{Source and licence.} Stability of shifts, interpolation, and crystalline measures. Version arXiv:2410.08682v2 (CC BY 4.0), distributed under the Creative Commons Attribution 4.0 International licence, as recorded in the arXiv licence metadata for this exact version and shown on the abstract page https://arxiv.org/abs/2410.08682v2. Full rights notice: licenses/arxiv-2410.08682-CC-BY-4.0.txt.\par\smallskip

\noindent\textit{Editorial scope.} Counted unit: the Claim whose source label is \texttt{c3} in the article (source0.tex lines 797--799, source label \texttt{c3}), delivered together with the article's own complete original proof of it (source0.tex lines 801--819, one \texttt{proof} environment of 19 lines). No mathematics is written, repaired or re-derived: statement and proof are the article's own text line for line. Required context retained uncounted: c2 (lines 775--779). Every definition, display and label of those lines is retained unchanged. Omitted from this package and not redistributed: 1--774, 780--796, 800--800, 820--978. Nothing inside a retained excerpt is omitted or altered: every retained body is delivered exactly as the article's lines are written, so no omission receipt of any kind accompanies this package. Citations: the retained excerpts carry no citation command at all, so no bibliography entry of the article is transcribed and no reference is redistributed. Reference adaptation: none was needed and none was made; every reference inside the delivered unit resolves inside it, so no unresolved reference or citation is produced. Printed slips kept literal: no printed word, symbol, hypothesis or number of the counted statement or of its proof was altered. Layout and class: the article's own class is \texttt{amsart} with packages including amssymb, amsmath, amsthm, graphicx, subcaption, stmaryrd, physics, inputenc, amsfonts, hyperref, xcolor; the wrapper is the lane's pinned \texttt{amsart} editorial wrapper at 10pt on A4, which restates the theorem-like environments the retained text needs and adds the article's own macro definitions that the retained text needs (\texttt{\textbackslash Ff}, \texttt{\textbackslash R}, \texttt{\textbackslash Zer}), with extra wrapper packages (bm, bbm, dsfont, xspace, cancel, mathtools). The delivered numbering is the unit's own. Assets: no retained span contains a figure environment, a graphics-inclusion command, a PDF-page inclusion, a table or a TikZ picture; no image file is copied into this package, the package declares no asset dependency and no image file is redistributed with it. Licence: version arXiv:2410.08682v2 of this article is distributed under the Creative Commons Attribution 4.0 International licence, as recorded in the arXiv licence metadata for this exact version and shown on the abstract page https://arxiv.org/abs/2410.08682v2. A full rights notice is delivered at licenses/arxiv-2410.08682-CC-BY-4.0.txt. Characters: every retained excerpt is pure ASCII, so the delivered engine, XeLaTeX with its default Latin Modern text font, reports no missing character.
\begin{claim}\label{c2}
    Let $I_0$ be an open interval and $t_0\in I_0$ such that $I_0\cap \Zer(\hat\Ff)=\{t_0\}.$
    Assume $\phi\in S(\R)$ vanishes outside $I_0$ and satisfies $\phi(t_0)=\phi'(t_0)=0.$
    Then $\langle\mu,\hat\phi\rangle=0.$
\end{claim}

\begin{claim}\label{c3}
    Let $I_0$ and  $t_0\in I_0$ satisfy the assumptions of Claim \ref{c2}. There is a constant $a_0$ such that the equation    \begin{equation}\label{eqe}\langle\mu,\hat\phi\rangle=a_0\phi(t_0)\end{equation} holds for every  $\phi\in S(\R)$ which vanishes outside $I_0$.
\end{claim}

\begin{proof}
We may assume that  $I_0=(-a,b)$ and $t_0=0$.
    Fix any functions  $\psi_j\in S(\R), j=1,2,$  which vanish outside $I_0$ and satisfy $\psi_1(0)=1,\psi'_1(0)=0, \psi_2(0)=0, \psi'_2(0)=1.$ Set
$$a_0:=\langle\mu,\hat\psi_1\rangle,\quad b_0:=\langle\mu,\hat\psi_2\rangle.$$For every function $\psi\in S(\R)$ which vanishes outside $I_0,$
the function $\phi(t):=\psi(t)-\psi(0)\psi_1(t)-\psi'(0)\psi_2(t)$ satisfies the assumptions of Claim \ref{c2}.
It follows that $\langle\mu,\hat\phi\rangle=0$ and so \begin{equation}\label{ab0}\langle \mu,\hat\psi \rangle =\psi(0)\langle\mu,\hat\psi_1\rangle+\psi'(0)\langle\mu,\hat\psi_2\rangle=a_0\psi(0)+b_0\psi'(0).\end{equation}It remains to prove that $b_0=0.$

Now, fix a small positive number $\epsilon$ and set  $\psi(t):=\psi_2(t/\epsilon)$. Then $\psi(0)=0,\psi'(0)=1/\epsilon$ and $\hat\psi(x)=\epsilon\hat\psi_2(\epsilon x).$ By \eqref{ab0},
$$|b_0|/\epsilon=|\langle\mu,\hat\psi\rangle|=\left|\epsilon\sum_{\gamma^\ast\in\Gamma^\ast}c_{\gamma^\ast}\hat\psi_2(\epsilon  \gamma^\ast)\right|\leq \|c_{\gamma^\ast}\|_\infty \epsilon\sum_{\gamma^\ast\in\Gamma^\ast}|\hat\psi_2(\epsilon  \gamma^\ast)|.$$

\begin{claim}
  There is a constant $C=C(\psi_2)$ such that   $$\epsilon\sum_{\gamma^\ast\in\Gamma^\ast}|\hat\psi_2(\epsilon  \gamma^\ast)|\leq C,\quad 0<\epsilon<1.$$
\end{claim}

This can be easily deduced from the assumptions that $\Gamma^\ast$ is separated and $\hat\psi_2\in S(\R)$.

Finally, we see that $|b_0|\leq C\epsilon,$ for some positive constant $C.$
 Since $\epsilon$ can be chosen arbitrarily small, we conclude that $b_0=0.$   The claim follows from \eqref{ab0}.
\end{proof}

\end{document}
